\section{Two correction terms, uniformly in a logarithmic shift}
\label{sec:corrections}

The longer expansion in the growth-rate section of Wikipedia can also be
obtained with an arithmetic remainder.  We give its coefficients explicitly
and bound the remainder, rather than invoke a general asymptotic-expansion
theorem.  This is the form of the expansion associated with
Moser--Wyman~\cite{moser-wyman}.

Put $r=W(n)$, $a=e^r$, and
\[
 U_j(x)=\sum_{k=0}^jS(j,k)x^k,\qquad
 U_j=U_j(r),\qquad v=U_2=r^2+r,\qquad b=av.
\]
For example $U_3=r^3+3r^2+r$ and
$U_4=r^4+6r^3+7r^2+r$.  Thus every $U_j$ below is an explicitly defined
polynomial, computable using the Stirling recurrence.
Define the following rational functions of $r$:
\begin{align*}
 P_0&=\frac{U_4}{8v^2}-\frac{5U_3^2}{24v^3},&
 P_1&=-\frac{U_3}{2v^2},&
 P_2&=-\frac1{2v},\\
 Q_0&=-\frac{U_6}{48v^3}+
       \frac{7U_3U_5}{48v^4}+\frac{35U_4^2}{384v^4}
       -\frac{35U_3^2U_4}{64v^5}+\frac{385U_3^4}{1152v^6},\\
 Q_1&=\frac{U_5}{8v^3}-\frac{35U_3U_4}{48v^4}
                        +\frac{35U_3^3}{48v^5},\\
 Q_2&=-\frac{5U_4}{16v^3}+\frac{35U_3^2}{48v^4},&
 Q_3&=\frac{5U_3}{12v^3},&
 Q_4&=\frac1{8v^2}.
\end{align*}
Write $P(r,h)=P_0+hP_1+h^2P_2$ and
$Q(r,h)=Q_0+hQ_1+h^2Q_2+h^3Q_3+h^4Q_4$.

\begin{theorem}[A uniform two-correction expansion]\label{thm:two-corrections}
Let $H\ge0$ be an integer, $n\ge3$, and $h\in\mathbb Z$ satisfy
$n+h\ge0$ and $|h|\le Hr$.  If $a\ge2^{12}$, then
\begin{equation}\label{eq:two-corrections}
 B_{n+h}=\frac{(n+h)!\,e^{a-1}}{r^{n+h}\sqrt{2\pi av}}
    \left(1+\frac{P(r,h)}a+\frac{Q(r,h)}{a^2}
                        +\mathcal E_{n,h}\right),
 \qquad
 |\mathcal E_{n,h}|\le\frac{2^{400}(H+1)^6}{a^3}.
\end{equation}
In particular this gives the displayed two-correction expansion uniformly for
$h=O(\log n)$, because $\log n=r+\log r\le2r$ for $r\ge1$.
\end{theorem}

\begin{proof}
Circular coefficient extraction at the same radius $r=W(n)$ gives
\[
 \frac{B_{n+h}r^{n+h}\sqrt{2\pi av}}{(n+h)!e^{a-1}}
 =\sqrt{\frac{av}{2\pi}}\int_{-\pi}^{\pi}
       \exp\bigl(\Phi(\theta)-ih\theta\bigr)\,d\theta,
\]
where $\Phi$ is the phase in \eqref{eq:growth-fourier}.  Set
$\varepsilon=a^{-1/2}$, $t=\sqrt{av}\,\theta$, and
$T=a^{1/12}$.  We analyze $|t|\le T$ and the two tails separately.

\paragraph{Taylor expansion of the phase.}
Induction on $j$ using the chain rule gives
\[
 \frac{d^j}{d\theta^j}e^{r e^{i\theta}}
   =i^j e^{r e^{i\theta}}U_j(r e^{i\theta}),
 \qquad
 \left|\frac{d^j}{d\theta^j}e^{r e^{i\theta}}\right|
   \le aU_j(r).
\]
The polynomial recurrence is $U_{j+1}(x)=x(U_j(x)+U_j'(x))$.
Since $r\ge1$ and $v\ge r^2$,
\begin{equation}\label{eq:touchard-derivative-bound}
 U_j(r)/v^{j/2}\le B_j.
\end{equation}
The explicit values $B_3=5,B_4=15,B_5=52,B_6=203,B_7=877,B_8=4140$
satisfy $B_j<j!$ for $3\le j\le8$.
Taylor's estimate through order seven therefore gives
\[
 \Phi(\theta)-ih\theta=-t^2/2+
       \sum_{j=1}^5\varepsilon^j A_j(t)+R(t),
 \qquad |R(t)|\le\varepsilon^6|t|^8,
\]
where
\begin{align*}
 A_1&=-i\left(\frac{ht}{\sqrt v}+
                         \frac{U_3t^3}{6v^{3/2}}\right),&
 A_2&=\frac{U_4t^4}{24v^2},\\
 A_3&=\frac{iU_5t^5}{120v^{5/2}},&
 A_4&=-\frac{U_6t^6}{720v^3},&
 A_5&=-\frac{iU_7t^7}{5040v^{7/2}}.
\end{align*}
In particular $|A_1|\le H|t|+|t|^3$ and
$|A_j|\le|t|^{j+2}$ for $2\le j\le5$.

\paragraph{An explicit sixth-order remainder for the exponential.}
Fix $t$ and put $Z(x)=\sum_{j=1}^5x^jA_j(t)$ for $0\le x\le\varepsilon$.
Since $a\ge2^{12}$, we have
$\varepsilon T=a^{-5/12}\le1/16$.  The only real parts in $Z$ come from
$A_2$ and $A_4$, and $A_4\le0$, so
\[
 \Re Z(x)\le x^2t^4\le t^2/256.
\]
For $1\le k\le5$,
\[
 |Z^{(k)}(x)|\le D(t):=600(H+1)(1+|t|^7).
\]
Indeed, every derivative coefficient of a monomial of degree at most five
is at most $5!=120$, $x\le1$, and
$H|t|+|t|^3+\cdots+|t|^7\le5(H+1)(1+|t|^7)$.
The sixth derivative of $e^Z$ is $e^Z$ times the complete Bell polynomial
in $Z',\ldots,Z^{(6)}$.  This follows by six applications of the product
rule, or from \eqref{eq:block-multiplicity}; its coefficient sum is
$B_6=203$, and every monomial has at most six factors.  Since $D\ge1$,
\[
 |(e^Z)^{(6)}|\le203D^6e^{t^2/256}.
\]
Let $C_j(t)$ be the coefficient of $x^j$ in the degree-five Taylor polynomial
of $e^{Z(x)}$ at zero.  Taylor's integral remainder yields
\[
 \left|e^{Z(\varepsilon)}-\sum_{j=0}^5\varepsilon^j C_j(t)\right|
 \le\varepsilon^6\frac{203}{720}D^6e^{t^2/256}.
\]
The additional phase remainder contributes at most
$\varepsilon^6|t|^8\exp(t^2/128)$, since
$|R|\le t^2(\varepsilon|t|)^6\le t^2/256$ and
$|e^R-1|\le|R|e^{|R|}$.  After multiplication by $e^{-t^2/2}$, the sum of
these two error bounds is at most
\begin{equation}\label{eq:mw-integrable-remainder}
 \varepsilon^6 2^{66}(H+1)^6(1+|t|^{42})e^{-t^2/4}.
\end{equation}
For example, use $600^6<2^{60}$,
$(1+|t|^7)^6\le2^5(1+|t|^{42})$, and
$|t|^8\le1+|t|^{42}$.
The normalized full-line integral of
$(1+|t|^{42})e^{-t^2/4}$ is
$\sqrt2(1+2^{21}41!!)<2^{149}$.
Thus the central integrated error is less than
$2^{220}(H+1)^6a^{-3}$.

\paragraph{Coefficient evaluation.}
The odd $C_j$ are imaginary odd polynomials, so their integrals on the
symmetric central interval vanish exactly.  The even coefficients needed are
\[
 C_0=1,\quad C_2=A_2+\tfrac12A_1^2,
 \quad
 C_4=A_4+A_1A_3+\tfrac12A_2^2+\tfrac12A_1^2A_2+\tfrac1{24}A_1^4.
\]
The full-line Gaussian averages of $C_2,C_4$ are $P(r,h),Q(r,h)$,
respectively.  Substituting the $A_j$ and using the even moments
$1,3,15,105,945,10395$ proves exactly the displayed coefficients.
These moments satisfy $\mu_{2j}=(2j-1)\mu_{2j-2}$ by integration by parts,
with $\mu_0=1$; all odd moments vanish.
This is a finite polynomial calculation, including all terms in $h$.

Replacing the central Gaussian polynomial integrals by full-line integrals
also has a bounded error.  The same derivative bound gives
\[
 |C_0|+|C_2|+|C_4|
 \le2^{66}(H+1)^6(1+|t|^{42}).
\]
On $|t|\ge T$, multiply the integrand by the upper bound
$|t|^{36}/T^{36}$ for $1$.  Since $T^{36}=a^3$, the resulting error is at most
\[
 \frac{2^{66}(H+1)^6}{a^3}
  \frac1{\sqrt{2\pi}}
       \int_{\mathbb R}(|t|^{36}+|t|^{78})e^{-t^2/2}\,dt
 <\frac{2^{350}(H+1)^6}{a^3}.
\]
Here $77!!\le78^{39}<2^{273}$ and the smaller moment is harmless.

\paragraph{The noncentral part of the original integral.}
The factor $e^{-ih\theta}$ has absolute value one, so the tail bounds
\eqref{eq:growth-middle-tail} and \eqref{eq:growth-far-tail} are unchanged.
The middle tail is at most $(24/T)e^{-T^2/48}$.  Using
$e^x\ge x^{18}/18!$ shows that this is at most
$24\cdot18!\cdot48^{18}a^{-3}<2^{170}a^{-3}$.
For the far tail, the elementary inequality $r\le e^{r/2}=\sqrt a$
implies $\sqrt b\le2a$ and $a/(6r)\ge\sqrt a/6$.
Consequently
\[
 3\sqrt b\,e^{-a/(6r)}\le6a e^{-\sqrt a/6}
 \le 8!\,6^9 a^{-3}<2^{44}a^{-3}.
\]
The inequality $r\le e^{r/2}$ follows already from
$1+r/2+r^2/8\ge r$.  The required central cutoff lies inside
$|\theta|\le1/r$, because $T=a^{1/12}\le\sqrt{av}/r$.

Adding the central error, the Gaussian polynomial tail, and the two original
tails gives less than $2^{400}(H+1)^6a^{-3}$.
This proves \eqref{eq:two-corrections}.
\end{proof}

Every constant in this proof is intentionally generous.  The uniform
remainder is the main point: it verifies the two orders of correction and
their dependence on $h$ without hiding an effectiveness hypothesis in an
$O$-symbol.  For an explicit accuracy request $1/q$, set
$C_H=2^{400}(H+1)^6$ and take
\[
 n\ge\max\{2^{24},(C_H(q+1))^2\}.
\]
The inequality $r\le\sqrt a$ gives $n=ar\le a^{3/2}$ and hence
$a\ge\sqrt n$ when $a\ge1$.  Thus this threshold ensures both
$a\ge2^{12}$ and $C_H/a^3\le1/q$ for $q\ge1$.
Fractional powers and the integrations can be replaced by the root and
Riemann-sum algorithms of Section~\ref{sec:foundations}; all bounds remain
primitive recursive in $H$ and $q$.
